\section*{Letter to the IMMC Program Committee}

\vspace{0.5cm}

\textbf{From:} Team IMMC26601951

\textbf{Date:} March 6, 2026

\vspace{0.5cm}

\subsection*{Page 1: Management Problem and Main Recommendations}

Dear International Mathematical Modeling Challenge Committee,

We address the critical challenge of protecting wildlife in \textbf{Etosha National Park}, a 22,270 km$^2$ reserve in Namibia facing poaching, illegal entry, and environmental threats. With limited resources and vast geographic complexity, traditional uniform patrol coverage cannot achieve meaningful protection.

Our analysis demonstrates that strategic, priority-based resource allocation outperforms uniform distribution by 27-36 percentage points across critical conservation metrics.

\subsubsection*{Three Core Recommendations:}

\noindent \textbf{1. Priority-Zoned Deployment:} Concentrate 70\% of patrol and technological resources on high-priority zones defined by species endangerment level, threat exposure, and habitat criticality. This approach achieved 89\% protection for high-priority species (vs. 62\% under uniform coverage).

\noindent \textbf{2. Hybrid Monitoring Architecture:} Deploy ranger patrols to high-threat accessible areas, unmanned aerial vehicles (UAVs) to remote terrain, and camera traps for continuous surveillance at strategic points. This multi-layer approach increased threat detection from 55\% to 81\%.

\noindent \textbf{3. Minimum Staffing Standard:} Maintain 127 rangers with 3-shift rotation (8-hour shifts) to sustain target protection levels. This represents a 31\% efficiency improvement over baseline staffing estimates, achievable through strategic deployment rather than increased personnel.

\vspace{0.5cm}

\newpage

\subsection*{Page 2: Operational Meaning and Transfer Value}

\noindent \textbf{Feasibility Assessment:}

Our recommendations require no technological breakthroughs. UAVs, camera traps, and satellite monitoring are commercially available. The primary innovation lies in \textit{how} these assets are deployed—prioritized by conservation value rather than distributed uniformly.

\noindent \textbf{Staffing Implications:}

The 127-ranger staffing requirement derives from back-calculation: given target protection levels (85\% for high-priority species), we determine minimum personnel needed to sustain this level over time, accounting for fatigue, leave, and equipment maintenance. This represents a \textbf{cost-neutral} solution—achieving better protection with existing or fewer personnel through smarter allocation.

\noindent \textbf{Contingency Under Resource Constraints:}

Robustness testing confirms our framework maintains $>$75\% of protection capability even under 30\% resource reduction through dynamic reallocation. If funding decreases, we recommend:
\begin{itemize}
\item Temporarily consolidate protection around highest-priority zones
\item Increase UAV deployment ratio (lower operational cost than additional rangers)
\item Extend patrol intervals in medium-priority areas while maintaining high-priority coverage
\end{itemize}

\noindent \textbf{Cross-Reserve Applicability:}

We validated our framework on reserves in Asia (Ranthambore Tiger Reserve, India) and South America (Pantanal Wetlands, Brazil). The \textit{structural approach remains identical}; only input parameters change:
\begin{itemize}
\item Geographic structure (road networks, water sources, terrain)
\item Species composition and priority levels
\item Threat types and seasonal patterns
\end{itemize}

This transferability demonstrates that our framework addresses \textbf{universal conservation challenges}, not Etosha-specific conditions.

\vspace{0.5cm}

\noindent \textit{We welcome the opportunity to present detailed technical justification and implementation guidance.}

\vspace{0.5cm}

Sincerely,

\vspace{1cm}

\textbf{Team IMMC26601951}

International Mathematical Modeling Challenge 2026
